\section{Problem Statement and Research Questions}\label{sec:intro_rq}

The preceding motivation established why reliable causer attribution
matters, both for BMW's quality management process and, more generally,
for industrial decision-making under uncertainty. Building on this
motivation, this section formalizes the resulting problem and derives the
research questions that guide the remainder of this thesis.

Despite the large amount of knowledge generated by quality incidents,
rework activities, inspections, supplier complaints, and problem-solving
reports, identifying the true causer of a newly detected defect remains a
largely manual task. Before corrective actions can be initiated, quality
engineers must determine whether a defect originates from a supplier, a
production process, an assembly step, or a logistics activity. Current
investigation practice relies on expert judgment and exact-match retrieval
of previous incidents: effective when an identical case has already been
encountered, but prone to failure for incomplete or previously unseen
defect scenarios. As a result, relevant evidence scattered across
multiple enterprise systems remains underutilized, leading to delayed
decisions, repeated investigations, inconsistent causer attribution, and
inefficient problem-solving.

Existing technical approaches only partially close this gap. Rule-based
systems are transparent but cannot generalize beyond predefined scenarios;
machine learning models detect complex patterns but operate as black
boxes; and knowledge graphs support structured retrieval but offer limited
support for reasoning and decision-making. No existing approach combines
historical enterprise knowledge with evidence-based reasoning and
explainability into a single decision-support system.

This thesis addresses this gap by investigating an evidence-centric
framework that integrates heterogeneous industrial knowledge into a
knowledge graph and combines evidence retrieval, multi-source evidence
aggregation, and explainable-AI techniques to generate transparent,
trustworthy recommendations for defect causer identification. This leads
to the following three research questions:

\begin{description}[%
  style=nextline,
  leftmargin=0pt,
  labelsep=0pt,
  font=\normalfont\bfseries,
  itemsep=\baselineskip,
  topsep=\smallskipamount
]
  \item[RQ1 -- Knowledge Retrieval and Evidence Discovery]
    How can heterogeneous industrial knowledge sources be integrated and
    retrieved through a knowledge graph to support the investigation of
    previously unseen quality incidents?

  \item[RQ2 -- Evidence Aggregation and Reasoning]
    How can evidence originating from multiple historical sources be
    aggregated, evaluated, and reasoned upon to support reliable defect
    causer identification and decision-making?

  \item[RQ3 -- Explainable Decision Support]
    How can generated recommendations be presented through transparent
    and auditable evidence-based explanations that support trust and
    acceptance among industrial domain experts?
\end{description}
